\documentclass[a4paper,11pt,oneside]{memoir}

% Castellano
\usepackage[spanish,es-tabla]{babel}
\selectlanguage{spanish}
\usepackage[utf8]{inputenc}
\usepackage{placeins} 
\usepackage{float}
\usepackage{eurosym}



\RequirePackage{booktabs}
\RequirePackage[table]{xcolor}
\RequirePackage{xtab}
\RequirePackage{multirow}

 \newcommand{\tabitem}{~\llap{\textbullet}~}

%Tablas
\usepackage[spanish,es-tabla]{babel}
\usepackage{longtable} % para tablas largas

%Simbolo del euro
\usepackage{eurosym} 

%Para codigo
\usepackage{listings}

%Directorios
\usepackage{forest}
\usetikzlibrary{shapes.geometric}
 \tikzset{/forest,
    symlink/.append style={
      opacity=.25,
      text opacity=.5,
      before drawing tree={
        {tikz+={\draw [thick, -{>[]}] (!#1.east) ++(4pt,-1.5pt) arc (215:120:5pt);}}
      },
    },
  }

 

% Links
\usepackage[colorlinks]{hyperref}
\hypersetup{
	colorlinks,
	linkcolor={white!60!black},
	citecolor={blue!50!black},
     urlcolor={blue!80!black}
}

% Ecuaciones
\usepackage{amsmath}

% Rutas de fichero / paquete
\newcommand{\ruta}[1]{{\sffamily #1}}

% Párrafos
\nonzeroparskip


% Landscape
\usepackage{pdflscape}


% Imagenes
\usepackage{graphicx}
\newcommand{\imagen}[2]{
	\begin{figure}[!h]
		\centering
		\includegraphics[width=0.9\textwidth]{#1}
		\caption{#2}\label{fig:#1}
	\end{figure}
	\FloatBarrier
}

\newcommand{\imagenflotante}[2]{
	\begin{figure}%[!h]
		\centering
		\includegraphics[width=0.9\textwidth]{#1}
		\caption{#2}\label{fig:#1}
	\end{figure}
}



% El comando \figura nos permite insertar figuras comodamente, y utilizando
% siempre el mismo formato. Los parametros son:
% 1 -> Porcentaje del ancho de página que ocupará la figura (de 0 a 1)
% 2 --> Fichero de la imagen
% 3 --> Texto a pie de imagen
% 4 --> Etiqueta (label) para referencias
% 5 --> Opciones que queramos pasarle al \includegraphics
% 6 --> Opciones de posicionamiento a pasarle a \begin{figure}
\newcommand{\figuraConPosicion}[6]{%
  \setlength{\anchoFloat}{#1\textwidth}%
  \addtolength{\anchoFloat}{-4\fboxsep}%
  \setlength{\anchoFigura}{\anchoFloat}%
  \begin{figure}[#6]
    \begin{center}%
      \Ovalbox{%
        \begin{minipage}{\anchoFloat}%
          \begin{center}%
            \includegraphics[width=\anchoFigura,#5]{#2}%
            \caption{#3}%
            \label{#4}%
          \end{center}%
        \end{minipage}
      }%
    \end{center}%
  \end{figure}%
}

%
% Comando para incluir imágenes en formato apaisado (sin marco).
\newcommand{\figuraApaisadaSinMarco}[5]{%
  \begin{figure}%
    \begin{center}%
    \includegraphics[angle=90,height=#1\textheight,#5]{#2}%
    \caption{#3}%
    \label{#4}%
    \end{center}%
  \end{figure}%
}
% Para las tablas
\newcommand{\otoprule}{\midrule [\heavyrulewidth]}
%
% Nuevo comando para tablas pequeñas (menos de una página).
\newcommand{\tablaSmall}[5]{%
 \begin{table}
  \begin{center}
   \rowcolors {2}{gray!35}{}
   \begin{tabular}{#2}
    \toprule
    #4
    \otoprule
    #5
    \bottomrule
   \end{tabular}
   \caption{#1}
   \label{tabla:#3}
  \end{center}
 \end{table}
}

%
% Nuevo comando para tablas pequeñas (menos de una página).
\newcommand{\tablaSmallSinColores}[5]{%
 \begin{table}[H]
  \begin{center}
   \begin{tabular}{#2}
    \toprule
    #4
    \otoprule
    #5
    \bottomrule
   \end{tabular}
   \caption{#1}
   \label{tabla:#3}
  \end{center}
 \end{table}
}



\newcommand{\tablaApaisadaSmall}[5]{%
\begin{landscape}
  \begin{table}
   \begin{center}
    \rowcolors {2}{gray!35}{}
    \begin{tabular}{#2}
     \toprule
     #4
     \otoprule
     #5
     \bottomrule
    \end{tabular}
    \caption{#1}
    \label{tabla:#3}
   \end{center}
  \end{table}
\end{landscape}
}

%
% Nuevo comando para tablas grandes con cabecera y filas alternas coloreadas en gris.
\newcommand{\tabla}[6]{%
  \begin{center}
    \tablefirsthead{
      \toprule
      #5
      \otoprule
    }
    \tablehead{
      \multicolumn{#3}{l}{\small\sl continúa desde la página anterior}\\
      \toprule
      #5
      \otoprule
    }
    \tabletail{
      \hline
      \multicolumn{#3}{r}{\small\sl continúa en la página siguiente}\\
    }
    \tablelasttail{
      \hline
    }
    \bottomcaption{#1}
    \rowcolors {2}{gray!35}{}
    \begin{xtabular}{#2}
      #6
      \bottomrule
    \end{xtabular}
    \label{tabla:#4}
  \end{center}
}

%
% Nuevo comando para tablas grandes con cabecera.
\newcommand{\tablaSinColores}[6]{%
  \begin{center}
    \tablefirsthead{
      \toprule
      #5
      \otoprule
    }
    \tablehead{
      \multicolumn{#3}{l}{\small\sl continúa desde la página anterior}\\
      \toprule
      #5
      \otoprule
    }
    \tabletail{
      \hline
      \multicolumn{#3}{r}{\small\sl continúa en la página siguiente}\\
    }
    \tablelasttail{
      \hline
    }
    \bottomcaption{#1}
    \begin{xtabular}{#2}
      #6
      \bottomrule
    \end{xtabular}
    \label{tabla:#4}
  \end{center}
}

%
% Nuevo comando para tablas grandes sin cabecera.
\newcommand{\tablaSinCabecera}[5]{%
  \begin{center}
    \tablefirsthead{
      \toprule
    }
    \tablehead{
      \multicolumn{#3}{l}{\small\sl continúa desde la página anterior}\\
      \hline
    }
    \tabletail{
      \hline
      \multicolumn{#3}{r}{\small\sl continúa en la página siguiente}\\
    }
    \tablelasttail{
      \hline
    }
    \bottomcaption{#1}
  \begin{xtabular}{#2}
    #5
   \bottomrule
  \end{xtabular}
  \label{tabla:#4}
  \end{center}
}



\definecolor{cgoLight}{HTML}{EEEEEE}
\definecolor{cgoExtralight}{HTML}{FFFFFF}

%
% Nuevo comando para tablas grandes sin cabecera.
\newcommand{\tablaSinCabeceraConBandas}[5]{%
  \begin{center}
    \tablefirsthead{
      \toprule
    }
    \tablehead{
      \multicolumn{#3}{l}{\small\sl continúa desde la página anterior}\\
      \hline
    }
    \tabletail{
      \hline
      \multicolumn{#3}{r}{\small\sl continúa en la página siguiente}\\
    }
    \tablelasttail{
      \hline
    }
    \bottomcaption{#1}
    \rowcolors[]{1}{cgoExtralight}{cgoLight}

  \begin{xtabular}{#2}
    #5
   \bottomrule
  \end{xtabular}
  \label{tabla:#4}
  \end{center}
}




\graphicspath{ {./img/} }

% Capítulos
\chapterstyle{bianchi}
\newcommand{\capitulo}[2]{
	\setcounter{chapter}{#1}
	\setcounter{section}{0}
	\chapter*{#2}
	\addcontentsline{toc}{chapter}{#2}
	\markboth{#2}{#2}
}

% Apéndices
\renewcommand{\appendixname}{Apéndice}
\renewcommand*\cftappendixname{\appendixname}

\newcommand{\apendice}[1]{
	%\renewcommand{\thechapter}{A}
	\chapter{#1}
}

\renewcommand*\cftappendixname{\appendixname\ }

% Formato de portada
\makeatletter
\usepackage{xcolor}
\newcommand{\tutor}[1]{\def\@tutor{#1}}
\newcommand{\course}[1]{\def\@course{#1}}
\definecolor{cpardoBox}{HTML}{E6E6FF}
\def\maketitle{
  \null
  \thispagestyle{empty}
  % Cabecera ----------------
\noindent\includegraphics[width=\textwidth]{cabecera}\vspace{1cm}%
  \vfill
  % Título proyecto y escudo informática ----------------
  \colorbox{cpardoBox}{%
    \begin{minipage}{.8\textwidth}
      \vspace{.5cm}\Large
      \begin{center}
      \textbf{TFG del Grado en Ingeniería Informática}\vspace{.6cm}\\
      \textbf{\LARGE\@title{}}
      \end{center}
      \vspace{.2cm}
    \end{minipage}

  }%
  \hfill\begin{minipage}{.20\textwidth}
    \includegraphics[width=\textwidth]{escudoInfor}
  \end{minipage}
  \vfill
  % Datos de alumno, curso y tutores ------------------
  \begin{center}%
  {%
    \noindent\LARGE
    Presentado por \@author{}\\ 
    en Universidad de Burgos --- \@date{}\\
    Tutor: \@tutor{}\\
  }%
  \end{center}%
  \null
  \cleardoublepage
  }
\makeatother


\newcommand{\nombre}{Adrian Aguado} %%% cambio de comando

% Datos de portada
\title{{\Huge barterAPP}\\[0.5cm]Manage your business time.}
\title{{DOCUMENTACIÓN TÉCNICA}}
\author{\nombre}
\tutor{Luis R.Izquierdo}
\date{\today}

\begin{document}

\maketitle


\cleardoublepage

%%%%%%%%%%%%%%%%%%%%%%%%%%%%%%%%%%%%%%%%%%%%%%%%%%%%%%%%%%%%%%%%%%%%%%%%%%%%%%%%%%%%%%%%

\frontmatter


\clearpage

% Indices
\tableofcontents

\clearpage

\listoffigures

\clearpage

%\listoftables

%\clearpage

\mainmatter

\appendix

\apendice{Plan de Proyecto Software}

\section{Introducción}\label{introduccion-plan}
En este capitulo se detalla la planificación del proyecto. Como gestor de tareas se comenzó utilizando \emph{Trello+Github} pero más tarde se pasó a utilizar \emph{Zenhub,} extensión de Google Chrome que permite integrar los \emph{boards} dentro del mismo repositorio de código alojado en \emph{github}. Se han dado más detalles en la memoria del proyecto.

Se han utilizado metodologías ágiles para el desarrollo del proyecto y de este modo, se ha realizado un desarrollo dividido en iteraciones. Terminada una iteración empezaba la siguiente y se agregaban a las tareas planeadas las que no habían sido completado de la iteracción precedente. Las iteraciones del proyecto estaban pensadas para durar una diez días aproximadamente. No obstante, hay alguna excepción en la que la iteración duró más tiempo. También existe alguna demora entre algún sprint debido a que tenía demasiada carga de trabajo de las asignaturas, trabajaba o estaba de viaje. 

La fase de planificación se puede dividir a su vez en:

\begin{itemize}
\tightlist
\item
  Planificación temporal.
\item
  Estudio de viabilidad.
\end{itemize} 

La primera parte me centro en la programación y desarrollo de la aplicación. Es decir elaboro un programa de tiempos con una serie de tareas a seguir para cumplimentar el proyecto.

La segunda parte se centra en el estudio de viabilidad. De la misma manera desde la segunda semana de marzo vengo realizando un plan de empresa con el \emph{programa Yuzz} por lo que ello me va a facilitar el estudio de viabilidad de mi proyecto. Se desarrollará tanto la viabilidad legal como también la económica. 


\section{Planificación temporal}\label{planificacion-temporal}
Desde inicio del proyecto se planteó utilizar una metodología ágil como
\emph{Scrum} para la gestión del proyecto. Aunque no se ha seguido al 100\% la
metodología al tratarse de un proyecto para la universidad, sí que se ha aplicado
en líneas generales una filosofía ágil y metódica. La diferencia fundamental radica en que esta metodología esta pensada para equipos y no para individuos.

A continuación se describen los diferentes \emph{sprints} que se han
realizado. Dentro de \emph{github} cada \emph{milestone} recibe el número del sprint asignado y dentro de cada uno de ellos existen diferentes tareas que describiré a continuación. A cada tarea le acompaña un número a la derecha el cuál es denominado \emph{story point}, de alguna manera sirve para realizar una estimación de lo que te va llevar completar esa determinada tarea. En mi caso concreto el \textbf{1} resulta ser el más bajo lo cuál indicaría que no más de dos o tres horas con cada tarea y el  \textbf{13} el más alto lo cuál significa varios días de trabajo.


\subsection{Sprint 0: 18/02/2017 - 28/02/2017}\label{sprint0}

Tareas principales:

\begin{itemize}
	\item Terminar formación.
	\item Aprender \emph{Sonarqube}.
	\item Inicio Back-End.
    \item Inicio Decidir base de datos a emplear.
\end{itemize}

\imagen{sprint0}{\footnotesize{Detalle sprint 0.}}

La primera vez que hice uso de \emph{ZenHub} ya llevaba algún tiempo formándome, de ahí el primer \emph{Issue} del \underline{Sprint 0}. Esta primera toma de contacto fue para comenzar a desarrollar el Back-End de la aplicación, además se consultaron varias fuentes para decidir que tipo de base de datos emplear.

\subsection{Sprint 1: 18/02/2017 - 28/02/2017}\label{sprint1}

Tareas principales:

\begin{itemize}
	\item Correción de errores en Back-End.
	\item Desarrollo Back-end.
	\item Errores en base de datos.
\end{itemize}

\imagen{sprint1}{\footnotesize{Detalle sprint 1.}}

El  \underline{Sprint 1} sirvió para continuar con el Back-end de la aplicación sin duda fue un de las partes más complicadas al pelearme con bases de datos con conceptos nuevos por lo que tuve numerosos bugs a la hora de guardar los usuarios en la base de datos.

\subsection{Sprint 2: 15/03/2017 - 31/03/2017}\label{sprint2}

Tareas principales:

\begin{itemize}
	\item Corrección de errores en Back-End.
	\item Inicio Front-End.
	\item Bug en base de datos.
	\item Comienzo a leer sobre la documentación.
\end{itemize}

\imagen{sprint2}{\footnotesize{Detalle sprint 2.}}

El \underline{Sprint 2} fue el momento donde una vez tenía un back-end sólido debía trasladarlo a la parte del usuario por lo que comencé a realizar la parte del front-end.

\subsection{Sprint 3: 07/04/2017 - 15/04/2017}\label{sprint3}

Tareas principales:

\begin{itemize}
	\item Subir documentación.
	\item Elección del calendario.
	\item Corrección en componentes.
\end{itemize}

\imagen{sprint3}{\footnotesize{Detalle sprint 3.}}

El \underline{Sprint 3} se centra en la parte del calendario sobre todo, además de algo de documentación y corregir los errores que he arrastrado del back-end.

\subsection{Sprint 4: 16/04/2017 - 22/04/2017}\label{sprint4}

Tareas principales:

\begin{itemize}
	\item Actualizar a Angular CLI.
	\item Heroku y MLab
	\item Documentación
\end{itemize}

\imagen{sprint4}{\footnotesize{Detalle sprint 4.}}

El \underline{Sprint 4} es más corto dado que requiere un menor tiempo en realizar las tareas.

\subsection{Sprint 5: 30/04/2017 - 07/05/2017}\label{sprint5}

Tareas principales:

\begin{itemize}
	\item Angular CLI.
	\item Bugs: base de datos 
	\item Bugs: calendario
\end{itemize}

\imagen{sprint5}{\footnotesize{Detalle sprint 5.}}

El \underline{Sprint 5} fue complicado debido a que cambiar a Angular CLI resulta más sencillo a la hora de desplegar en servidor pero hay que saber como funciona realmente los proyectos en Angular CLI.

\subsection{Sprint 6: 09/05/2017 - 16/05/2017}\label{sprint6}

Tareas principales:

\begin{itemize}
	
		\item Error en turnos diarios, se decide implementar turnos 24 horas para comprobar si el algoritmo intercambia los turnos.
	\item Anexos en la documentación.
\end{itemize}

\imagen{sprint6}{\footnotesize{Detalle sprint 6.}}

El \underline{Sprint 6} se centra en que el algoritmo intercambie los turnos de manera eficiente. 

\subsection{Sprint 7: 24/05/2017 - 31/05/2017}\label{sprint7}

Tareas principales:

\begin{itemize}

		\item Tras consultar a Luis he obviado un importante elemento y es necesario refactorizar el algoritmo.
	\item Se incluye una parte para tests pero no me da tiempo.
\end{itemize}

\imagen{sprint7}{\footnotesize{Detalle sprint 7.}}

El \underline{Sprint 7} se centra en en conseguir que el algoritmo funcione para todos los tipos de turnos.

\subsection{Sprint 8: 01/06/2017 - 10/06/2017}\label{sprint8}

Tareas principales:

\begin{itemize}
	\item Completar últimos componentes
	\item Terminar documentación.
	\item Bugs: Travis. 
	\item Bugs: Heroku y fallo en el algoritmo.
	\item Pruebas de calidad del código (Cómo \emph{SonarQube} no da soporte a Typescript decidí prescindir de esta herramienta)
\end{itemize}

\imagen{sprint8}{\footnotesize{Detalle sprint 8.}}

El \underline{Sprint 8} se centra en corregir errores, uno de los cuales no consigo solventar, y realizar pruebas en la calidad del código. 

\subsection{Sprint 9: 10/06/2017 - 20/06/2017}\label{sprint9}

Tareas principales:

\begin{itemize}
	\item Refactorización.
	\item App móvil: scafolding básico y creación de componentes.
	\item Envío a Luis de la documentación para posibles comentarios.
	\item Error en el servidor.
	
\end{itemize}

\imagen{sprint9}{\footnotesize{Detalle sprint 9.}}

El \underline{Sprint 9} trata de refactorizar la app para que la calidad del código aumente. Se consigue en un tanto por ciento considerable. 

\subsection{Sprint 10: 20/06/2017 - 30/06/2017}\label{sprint10}

Tareas principales:

\begin{itemize}
	\item 
\end{itemize}

%\imagen{sprint10}{\footnotesize{Detalle sprint 10.}}

El \underline{Sprint 10} 

\subsection{Sprint 11: 30/06/2017 - 10/07/2017}\label{sprint11}

Tareas principales:

\begin{itemize}
	\item 
\end{itemize}

%\imagen{sprint10}{\footnotesize{Detalle sprint 11.}}

El \underline{Sprint 11} 


\section{Estudio de viabilidad}\label{estudio-viabilidad}
En esta sección se lleva a cabo un estudio para comprobar la viabilidad del proyecto realizado. Paralelamente al desarrollo de la aplicación, como ya se ha nombrado anteriormente, el proyecto formó parte del \emph{programa YUZZ } para jóvenes emprendedores en el que durante cinco meses realicé un plan de empresa completo. Se detalla por tanto en un documento extra que adjuntaré un estudio de viabilidad exhaustivo  y muy completo en el que se incluyen entre otras cosas: plan de marketing, plan de financiación, estudio de viabilidad o plan de puesta en marcha del negocio a cinco años vista. Así como también encuestas, logotipado o un \emph{landing page} necesaria para la presentación del proyecto ante el jurado. 

Por lo tanto en esta sección voy a realizar un resumen del documento descrito en el párrafo anterior en el que como conclusión definitiva tendremos un boceto de lo que supondría transformar un proyecto fin de carrera y que pase a formar parte del mercado. Así mismo voy a intentar adaptarlo a las condiciones que se exigen en el proyecto dado que el plan de empresa adjunto como extra es un estudio de viabilidad íntegro de aquí a cinco años por lo que resulta ser más extenso y detallado. Se intentará por tanto aquí realizar una estimación.

\subsection{Viabilidad económica}\label{viabilidad-economica}

La viabilidad económica es la parte donde lograremos detectar si el proyecto es o no rentable económicamente hablando.

\subsubsection{Análisis de costes}\label{costes}

Lo subdividiremos en diferentes tipos de costes.

\begin{description}
	\item[Coste de personal] Se considerará que el proyecto ha sido desarrollado en un periodo de cinco meses.  Considerando que se ha trabajado unas 6 horas al día cada semana, y que el programador, que en este caso es una sola persona, ha percibido un sueldo de 13 \euro /hora, el coste del personal por lo tanto se resume en la siguiente tabla:
	
\begin{table}[htbp]
\begin{center}
\begin{tabular}{|l|l|}
\hline
 & Total \\
\hline \hline
13 \euro /hora * 6 horas/día &   78 \euro /día \\ \hline
78 \euro /día * 5 dias/semana &   390 \euro /semana \\ \hline
390 \euro /semana * 4 semanas/mes &   1560 \euro /mes \\ \hline
\textbf{Coste total salario} &   7800 \euro /5 meses \\ \hline
\end{tabular}
\caption{Tabla salarios.}
\label{tabla:salarios}
\end{center}
\end{table}
	
	\item[Coste de seguridad social] al coste de personal hay que sumar el coste de seguridad social. 
	
	
	\begin{table}[htbp]
\begin{center}
\begin{tabular}{|l|l|}
\hline
 \emph{Seguridad social} &  \\
\hline \hline
 & \emph{23,60} \% en contigencias comunes \\ \hline
&   \emph{7,70} \% en concepto de desempleo \\ \hline
 &   \emph{0,10} \% en concepto de formación personal. \\ \hline
\textbf{Gastos Seguridad Social}  &   7800 \euro  * 31,40 \% =  2449,20 \euro \\ \hline
\textbf{Gastos personal}  &  7800 \euro  + 2449,20 \euro  = 10249,2 \euro  \\ \hline
\end{tabular}
\caption{Tabla seguridad social.}
\label{tabla:ssocial}
\end{center}
\end{table}
	
	\item[Coste de software] Todo el software utilizado para este proyecto es gratuito y de libre distribución. 
	
	\item[Coste de Hardware] En este apartado se analizan los costes derivados de la compra de los materiales que son necesarios para llevar a cabo el proyecto. Según la ley de Impuesto de Sociedades, la vida media del inmovilizado  hardware oscila entre 8 y 3 a\~nos. Yo consideraré 3 años. Más detalles en la siguiente tabla.  
	
\begin{table}[htbp]
\begin{center}
\begin{tabular}{|l|l|}
\hline
 \emph{Hardware utilizado}&  \\
\hline \hline
 Portátil MacBook Pro 15'' & 1650 \euro  \\ \hline
Monitor externo Samsung 22'' &    107'99 \euro\\ \hline
\textbf{Gastos Hardware}  &   1650 \euro + 107,99  =  1757,99 \euro \\ \hline
\textbf{Coste amortización}  &  1757,99 \euro  / 36 meses = 48,8 \euro / mes  \\ \hline
\end{tabular}
\caption{Tabla costes hardware.}
\label{tabla:hardware}
\end{center}
\end{table}

\begin{table}[htbp]
\begin{center}
\begin{tabular}{|l|l|}
\hline
\textbf{Coste final hardware} (5 meses)  &  48,8 \euro / mes  *  5 meses = 244 \euro   \\ \hline
\end{tabular}
\caption{Coste final hardware.}
\label{tabla:costefinalhardware}
\end{center}
\end{table}


\subsection{Coste final desarrollo del proyecto}\label{coste-final}

El coste final del desarrollo del proyecto, con una duración estimada de cinco meses es el siguiente.  

\begin{table}[htb]
\begin{center}
\begin{tabular}{|l|l|}
\hline
\textbf{Coste total final del desarrollo proyecto}  &  244 \euro +  10249,2 \euro = 10493 \euro   \\ \hline
\end{tabular}
\caption{Coste proyecto final..}
\label{tabla:costefinal}
\end{center}
\end{table}

\end{description}

\subsection{Viabilidad legal}\label{viabilidad-legal}
La viabilidad legal se centra principalmente en el estudio de las licencias software utilizadas. Realizaremos una tabla resumen sobre las partes más importantes de la aplicación y la licencia que emplean. Realmente son muchos los módulos que se han empleado para el desarrollo pero la mayoría tienen la misma licencia por lo que no he nombrado todos simplemente los que más importantes parecen. 

\begin{table}[H]
\begin{center}
\begin{tabular}{|l|l|}
\hline
Dependencias & Licencia \\
\hline \hline
\emph{angular-calendar} & \href{https://opensource.org/licenses/MIT}{MIT}\\ \hline
\emph{@angular/cli} & \href{https://opensource.org/licenses/MIT}{MIT}\\ \hline
\emph{@types/node} & \href{https://opensource.org/licenses/MIT}{MIT}\\ \hline
\emph{primeng} & \href{https://opensource.org/licenses/MIT}{MIT}\\ \hline
\emph{tslint} & \href{https://opensource.org/licenses/Apache-2.0}{Apache-2.0}\\ \hline
\emph{typescript} & \href{https://opensource.org/licenses/Apache-2.0}{Apache-2.0}\\ \hline
\emph{font-awesome} & \href{https://opensource.org/licenses/OFL-1.1}{OFL-1.1}\\ \hline
\emph{mongoose} & \href{https://opensource.org/licenses/MIT}{MIT}\\ \hline
\emph{express} & \href{https://opensource.org/licenses/MIT}{MIT}\\ \hline
\emph{bootstrap} & \href{https://opensource.org/licenses/MIT}{MIT}\\ \hline

\end{tabular}
\caption{Tabla resumen-licencias.}
\label{tabla:licencias}
\end{center}
\end{table}



Una licencia software es un contrato entre el autor o titular de los derechos de explotación o distribución y el usuario consumidor, usuario profesional o empresa, para la utilización del software cumpliendo una serie de términos y condiciones establecidas dentro de sus cláusulas. Todo el software que empleo es su amplia mayoría es libre ya que la mayoría de código viene o bien de desarrolladores \emph{amateur} o es libre desde su creación. Emplean por tanto las licencias descritas en la tabla anterior por lo que diferencia entre ellas varían en términos como el nombramiento del autor o la garantía , son licencia comunes en software libre \cite{githublicense} , además github provee una página para ayudarte en la elección de tu licencia \citep{githubchoose}. La menos permisiva de entre las descritas es \emph{MIT} \citep{mit}  y la más es la licencia del  \emph{Apache-2.0} \cite{apache}.

\imagen{licencias}{\footnotesize{Tanto por ciento licencias utilizadas.}}

Tal y cómo he nombrado anteriormente, y el tutor así lo deseaba, el software puede ser modificado, ayudando así a que crezca en un futuro
por lo que he elegido \emph{Apache-2.0}. 


Para la documentación he elegido una licencia Creative commons, en concreto se ha elegido la \emph{Attribution-NonCommercial 4.0 International (CC BY-NC 4.0)}


\imagen{creativecommons}{\footnotesize{Licencia Creative Commons.}}




\apendice{Especificación de Requisitos}

\section{Introducción}\label{introduccion-requisitos}

Este anexo recoge la especificación de requisitos que define el comportamiento del sistema desarrollado. El objetivo principal de la Especificación de Requisitos del Sistema (\emph{ERS}) es servir como medio de comunicación entre clientes, usuarios, ingenieros de requisitos y desarrolladores.

La ERS es correcta si y sólo si todo requisito que figura en ella refleja alguna necesidad
real. La corrección de la ERS implica que el sistema implementado será el sistema
deseado. Se han seguido las recomendaciones del estándar IEEE 830 según la última versión del estándar citado. Las características deseables para una especificación de requisitos son:


\begin{enumerate}
	\item No ambigua.
	\item Completa
	\item  Verificable
	\item  Consistente
	\item  Clasificada
	\item  Modificable
	\item  Explorable
	\item  Utilizable durante las tareas de mantenimiento y uso
\end{enumerate}


\section{Objetivos generales}\label{objetivos-generales}
Los objetivos generales que se perseguían con el proyecto han sido:
\begin{itemize}
	\item Desarrollar una aplicación web que permita intercambiar turnos.
	\item Desarrollar una aplicación híbrida que adapte la anterior aplicación web a cualquier dispositivo.
\end{itemize}


\section{Catálogo de requisitos}\label{catalogo-requisitos}
A continuación, se enumeran los requisitos específicos:

En esta sección se especificarán los requisitos del sistema, diferenciando los requisitos funcionales, o sea el comportamiento del sistema, de los requisitos no funcionales, que describen características de funcionamiento.
\subsection{Requisitos funcionales}\label{r-funcionales}
\begin{description}
    \item[RF-1 Gestión de usuarios:] Los usuarios son una parte fundamental.
    \begin{itemize}
         \item \textbf{RF-1.1:} Registrar usuarios
         \item \textbf{RF-1.2:} Logear usuarios
         \item \textbf{RF-1.3:} Deslogearse
    \end{itemize}
   \item[RF-2. Gestión de calendario:] Funciona como un calendario normal.
    \begin{itemize}
         \item \textbf{RF-2.1:} Añadir turnos.
         \item \textbf{RF-2.2:} Eliminar turnos.
         \item \textbf{RF-2.3:} Editar turnos.
         \item \textbf{RF-2.4:} Visualizar turnos.
    \end{itemize}
 	   \item[RF-3. Gestión de cambios:] Parte de intercambio de turnos.
    \begin{itemize}
         \item \textbf{RF-3.1:} Enviar petición.
         \item \textbf{RF-3.2:} Aceptar petición.
         \item \textbf{RF-3.3:} Rechazar petición. 
         \item \textbf{RF-3.4:} Aceptar intercambio.
         \item \textbf{RF-3.5:} Rechazar intercambio.
         \item \textbf{RF-3.6:} Visualizar detalles.
           \begin{itemize}
             \item \textbf{RF-3.6.1:} Visualizar turnos.
             \item \textbf{RF-3.6.2:} Visualizar pendientes.
             \item \textbf{RF-3.6.3:} Visualizar aceptados.
             \item \textbf{RF-3.6.4:} Visualizar rechazados.
           \end{itemize}
         \item \textbf{RF-3.7:} Visualizar calendario.
    \end{itemize}
       
    
    
   \item[RF-4. Información:] Información de como usar la aplicación
\end{description}

\subsection{Requisitos no funcionales}\label{rnofuncionales}
\begin{description}
    \item [RNF-1 Seguridad] se persigue preservar la seguridad impidiendo acceder a la aplicación sin logearse. O también añadiendo un \emph{hash} para el almacenamiento de las contraseñas de los usuarios, preservando así la seguridad de los mismos.
    \item [RNF.2 Escalabilidad:] la aplicación debe estar preparada para actuar con un gran número de usuarios.
    \item [RNF.3 Soporte:]  la aplicación debe ser soportable por diferentes navegadores (Restricción: lamentablemente \emph{Angular 4} no es soportado actualmente por todos los navegadores del mercado)
    \item [RNF.4 Usabilidad:]  la aplicación debe ser fácil de usar.
\end{description}


\section{Especificación de requisitos}\label{requisitos}

En esta sección se explicará el diagrama de casos de uso y se desarrollará cada uno de los requisitos en función del esquema.

\subsection{Diagrama casos de uso}

%Imagen casos de uso
\imagen{casodeusoschema}{\footnotesize{Casos de uso. Fuente: Elaboración propia.}}

\subsection{Actores}
El actor es el usuario de la aplicación tenemos que tener en cuenta que para que se de un cambio de turno deben de existir dos actores sino sería imposible el intercambio.

\subsection{Casos de uso}

%%%%%%%%%%%%%%%%%%%%%%%%%%%
%CU-1: Gestión usuarios
\begin{table}[H]
\begin{tabular}{llr}  
\toprule
\begin{minipage}[b]{0.24\columnwidth}\raggedright\strut
\textbf{CU-01}\strut
\end{minipage} & \begin{minipage}[b]{0.72\columnwidth}\raggedright\strut
\textbf{Gestión usuarios}\strut
\end{minipage}\tabularnewline
\cmidrule(r){1-2}
\textbf{Versión}       & 1.0           \\
\textbf{Autor}       & Adrián  Aguado    \\
\textbf{Requisitos asociados}       & RF-1.1, RF-1.2, RF-1.3 \\ 
\textbf{Descripción} & Permite registrarse, logearse al usuario o salir.\\
\textbf{Precondición} & La página está cargada, \\
& base de datos funcionando.       \\
\textbf{Acciones posibles} & 1. El usuario se registra. \\
& 2. El usuario  se logea. \\
& 3. El usuario interactúa.\\
& 4. El usuario sale.        \\
\textbf{Postcondición} & Se añade el usuario \\
& a la base de datos.     \\
\textbf{Excepciones} &  \\
\textbf{Frecuencia} & Alta            \\
\textbf{Importancia} & Alta            \\
\textbf{Comentarios } &       \\
\bottomrule
\end{tabular}
\caption{CU-1 Gestión usuarios} 
\end{table}

%CU-02: Registro
\begin{table}[H]
\begin{tabular}{llr}  
\toprule
\begin{minipage}[b]{0.24\columnwidth}\raggedright\strut
\textbf{CU-02}\strut
\end{minipage} & \begin{minipage}[b]{0.72\columnwidth}\raggedright\strut
\textbf{Registro}\strut
\end{minipage}\tabularnewline
\cmidrule(r){1-2}
\textbf{Versión}       & 1.0           \\
\textbf{Autor}       & Adrián  Aguado    \\
\textbf{Requisitos asociados}       & RF-1.1 \\ 
\textbf{Descripción} & Permite registrarse.\\
\textbf{Precondición} & La página está cargada  \\
& en la pestaña adecuada.     \\
\textbf{Acciones} & 1. El usuario introduce los datos, \\
& a saber: \\
& A. Nombre \\
& B. Apellidos\\
& B. Nombre usuario\\
& C. E-mail\\
& D. Nombre empresa\\
& E. Tipo de turno (a elegir)\\
& F. Contraseña\\
\textbf{Postcondición} & Se añade el usuario \\
& a la base de datos.     \\
& Redirección a log-in.  \\
\textbf{Excepciones} &  Si no se introducen todos los campos  \\
& obligatorios se notifica. También si el usuario \\
& ya existe o si todo ha ido correctamente.    \\
\textbf{Frecuencia} & Alta            \\
\textbf{Importancia} & Alta            \\
\textbf{Comentarios } &       \\
\bottomrule
\end{tabular}
\caption{CU-02 Registro} 
\end{table}

%CU-03: Log-in
\begin{table}[H]
\begin{tabular}{llr}  
\toprule
\begin{minipage}[b]{0.24\columnwidth}\raggedright\strut
\textbf{CU-03}\strut
\end{minipage} & \begin{minipage}[b]{0.72\columnwidth}\raggedright\strut
\textbf{Log-in}\strut
\end{minipage}\tabularnewline
\cmidrule(r){1-2}
\textbf{Versión}       & 1.0           \\
\textbf{Autor}       & Adrián  Aguado    \\
\textbf{Requisitos asociados}       & RF-1.2 \\ 
\textbf{Descripción} & Permite logearse al usuario\\
\textbf{Precondición} & La página está cargada \\
& en la pestaña adecuada.     \\
& Base de datos funcionando.       \\
\textbf{Acciones} & 1. El usuario se \emph{logea}\\
& para ellos introduce:\\
& A. Correo electrónico \\
& B. Contraseña\\
& 2. Se cargan los turnos en el calendario.\\
\textbf{Postcondición} & El usuario entra en la  \\
& aplicación.    \\
\textbf{Excepciones} &  Si la contraseña no es correcta, aviso. \\
&  Si el correo no es correcto, aviso. \\
\textbf{Frecuencia} & Alta            \\
\textbf{Importancia} & Alta            \\
\textbf{Comentarios } &  Si fuera el primer \emph{log-in} el calendario   \\
& el calendario estará vacío.\\
\bottomrule
\end{tabular}
\caption{CU-03 Log-in} 
\end{table}

%CU-04: Log out
\begin{table}[H]
\begin{tabular}{llr}  
\toprule
\begin{minipage}[b]{0.24\columnwidth}\raggedright\strut
\textbf{CU-04}\strut
\end{minipage} & \begin{minipage}[b]{0.72\columnwidth}\raggedright\strut
\textbf{Log-out}\strut
\end{minipage}\tabularnewline
\cmidrule(r){1-2}
\textbf{Versión}       & 1.0           \\
\textbf{Autor}       & Adrián  Aguado    \\
\textbf{Requisitos asociados}       & RF-1.3 \\ 
\textbf{Descripción} & Permite salir de la aplicación.\\
\textbf{Precondición} & El usuario está dentro de la \\
& aplicación.       \\
\textbf{Acciones} & 1. Presionar el botón \emph{log out}.\\
   \\
\textbf{Postcondición} & Sale de la aplicación.  \\
\textbf{Excepciones} &     \\
\textbf{Frecuencia} & Alta            \\
\textbf{Importancia} & Alta            \\
\textbf{Comentarios } &      \\
\bottomrule
\end{tabular}
\caption{CU-04 Log out} 
\end{table}

Nota: El editar usuarios y eliminar usuarios no esta implementado de cara al usuario. Está implementado para realizar en una segunda fase  pero no se encuentra disponible en pantalla para realizarse directamente por parte del usuario.

%%%%%%%%%%%%%%%%%%%%%%%%%%%
%CU-05: Gestión Calendario
\begin{table}[H]
\begin{tabular}{llr}  
\toprule
\begin{minipage}[b]{0.24\columnwidth}\raggedright\strut
\textbf{CU-05}\strut
\end{minipage} & \begin{minipage}[b]{0.72\columnwidth}\raggedright\strut
\textbf{Gestión calendario}\strut
\end{minipage}\tabularnewline
\cmidrule(r){1-2}
\textbf{Versión}       & 1.0           \\
\textbf{Autor}       & Adrián  Aguado    \\
\textbf{Requisitos asociados}       & RF-2.1, RF-2.2, RF-2.3, RF-2.4  \\ 
\textbf{Descripción} & Gestión del calendario.\\
\textbf{Precondición} & El usuario está dentro de la \\
& aplicación.       \\
\textbf{Acciones} & 1. Añadir turnos. \\
& 2. Eliminar turnos. \\
& 3. Editar turnos. \\
& 4. Visualizar turnos. \\
\textbf{Postcondición} & Turno deseado. \\
\textbf{Excepciones} &     \\
\textbf{Frecuencia} & Media          \\
\textbf{Importancia} & Alta            \\
\textbf{Comentarios } &      \\
\bottomrule
\end{tabular}
\caption{CU-05 Gestión Calendario} 
\end{table}

%CU-06: Añadir turno
\begin{table}[H]
\begin{tabular}{llr}  
\toprule
\begin{minipage}[b]{0.24\columnwidth}\raggedright\strut
\textbf{CU-06}\strut
\end{minipage} & \begin{minipage}[b]{0.72\columnwidth}\raggedright\strut
\textbf{Añadir turno}\strut
\end{minipage}\tabularnewline
\cmidrule(r){1-2}
\textbf{Versión}       & 1.0           \\
\textbf{Autor}       & Adrián  Aguado    \\
\textbf{Requisitos asociados}       & RF-2.1 \\ 
\textbf{Descripción} & Añadir un turno determinado. \\
\textbf{Precondición} & El usuario está dentro de la \\
& aplicación.       \\
\textbf{Acciones} & 1. Presionar  botón \emph{nuevo turno}. \\
& 2. Seleccionar título. \\
& 3. Seleccionar tipo. \\
& 4 .Seleccionar fecha.\\
\textbf{Postcondiciones} & Turno deseado en calendario. \\
& Turno añadido en la base de \\
&  datos asociado al usuario \\
&  que lo ha añadido.  \\
&  Prioridad normal.  \\
\textbf{Excepciones} &  No es posible más de un turno al día.   \\
\textbf{Frecuencia} & Media          \\
\textbf{Importancia} & Alta            \\
\textbf{Comentarios } &     \\
\bottomrule
\end{tabular}
\caption{CU-06 Añadir turno} 
\end{table}




%CU-07: Eliminar turno
\begin{table}[H]
\begin{tabular}{llr}  
\toprule
\begin{minipage}[b]{0.24\columnwidth}\raggedright\strut
\textbf{CU-07}\strut
\end{minipage} & \begin{minipage}[b]{0.72\columnwidth}\raggedright\strut
\textbf{Eliminar turno}\strut
\end{minipage}\tabularnewline
\cmidrule(r){1-2}
\textbf{Versión}       & 1.0           \\
\textbf{Autor}       & Adrián  Aguado    \\
\textbf{Requisitos asociados}       & RF-2.3 \\ 
\textbf{Descripción} & Eliminar un turno.\\
\textbf{Precondición} & El usuario está dentro de la \\
& aplicación.       \\
\textbf{Acciones} & 1. Presionar botón \emph{eliminar turno} \\
& 2. Turno eliminado. \\
\textbf{Postcondiciones} & Turno eliminado del calendario.\\
& Turno eliminado de la base de datos. \\
\textbf{Excepciones} &     \\
\textbf{Frecuencia} & Media          \\
\textbf{Importancia} & Alta            \\
\textbf{Comentarios } &      \\
\bottomrule
\end{tabular}
\caption{CU-07 Eliminar turno} 
\end{table}

%CU-08: Editar turno
\begin{table}[H]
\begin{tabular}{llr}  
\toprule
\begin{minipage}[b]{0.24\columnwidth}\raggedright\strut
\textbf{CU-08}\strut
\end{minipage} & \begin{minipage}[b]{0.72\columnwidth}\raggedright\strut
\textbf{Editar turno}\strut
\end{minipage}\tabularnewline
\cmidrule(r){1-2}
\textbf{Versión}       & 1.0           \\
\textbf{Autor}       & Adrián  Aguado    \\
\textbf{Requisitos asociados}       & RF-2.3 \\ 
\textbf{Descripción} & Editar un turno.\\
\textbf{Precondición} & El usuario está dentro de la \\
& aplicación. \\
\textbf{Acciones} & 1.El usuario visualiza el turno  \\
& en la pestaña  \emph{normal}. \\
& 2. El usuario puede modificar :      \\
& A.Titulo \\
& B.Tipo \\
& C.Fecha\\
\textbf{Postcondición} & Turno deseado editado.\\
\textbf{Excepciones} &     \\
\textbf{Frecuencia} & Media          \\
\textbf{Importancia} & Alta            \\
\textbf{Comentarios } &    El usuario puede modificar  \\
&  un turno siempre que  \\
&  así lo desee.  \\
\bottomrule
\end{tabular}
\caption{CU-08 Editar turno} 
\end{table}

%CU-09: Visualizar turnos
\begin{table}[H]
\begin{tabular}{llr}  
\toprule
\begin{minipage}[b]{0.24\columnwidth}\raggedright\strut
\textbf{CU-09}\strut
\end{minipage} & \begin{minipage}[b]{0.72\columnwidth}\raggedright\strut
\textbf{ Visualizar turnos}\strut
\end{minipage}\tabularnewline
\cmidrule(r){1-2}
\textbf{Versión}       & 1.0           \\
\textbf{Autor}       & Adrián  Aguado    \\
\textbf{Requisitos asociados}       & RF-2.4 \\ 
\textbf{Descripción} & Ver de un simple vistazo \\
&todos los turnos. \\
\textbf{Precondición} & El usuario está dentro de la \\
& aplicación.       \\
\textbf{Acciones} & 1. Ver turnos. \\
\textbf{Postcondición} &   \\
\textbf{Excepciones} &     \\
\textbf{Frecuencia} & Media          \\
\textbf{Importancia} & Alta            \\
\textbf{Comentarios } &      \\
\bottomrule
\end{tabular}
\caption{CU-09 Visualizar turnos} 
\end{table}

%%%%%%%%%%%%%%%%%%%%%%%%%%%
%CU-10: Gestión cambios
\begin{table}[H]
\begin{tabular}{llr}  
\toprule
\begin{minipage}[b]{0.24\columnwidth}\raggedright\strut
\textbf{CU-10}\strut
\end{minipage} & \begin{minipage}[b]{0.72\columnwidth}\raggedright\strut
\textbf{Gestión cambios}\strut
\end{minipage}\tabularnewline
\cmidrule(r){1-2}
\textbf{Versión}       & 1.0           \\
\textbf{Autor}       & Adrián  Aguado    \\
\textbf{Requisitos asociados}       & RF-3.1, RF-3.2, RF-3.3, RF-3.4  \\ 
& RF-3.5, RF-3.6,RF-3.6.1, RF-3.6.2 \\ 
& RF-3.6.3, RF-3.6.4,  RF-3.7 \\ 
\textbf{Descripción} & Gestión de los cambios \\
& entre usuarios. \\
\textbf{Precondición}  & El usuario está dentro de la \\
& aplicación.      \\
& Hay más de un usuario. \\
\textbf{Acciones} & 1. Enviar petición de cambio. \\
& 2. Aceptar petición de cambio. \\
& 3. Rechazar petición de cambio. \\
& 4. Aceptar intercambio. \\
& 5. Rechazar intercambio.  \\
& 6. Visualizar detalles.  \\
& 7. Visualizar calendario. \\
\textbf{Postcondición} & Cambio turno. \\
\textbf{Excepciones} &   Para diversos casos  \\
 &   se muestran avisos.  \\
\textbf{Frecuencia} & Media          \\
\textbf{Importancia} & Alta            \\
\textbf{Comentarios } & Se controlan que la       \\
& petición se envíe una sola vez. \\
& Se muestra la lista de usuarios con \\
& turnos libres asociada a un \\
& día concreto \\
\bottomrule
\end{tabular}
\caption{CU-10 Gestión Cambios} 
\end{table}

%CU-11: Enviar petición
\begin{table}[H]
\begin{tabular}{llr}  
\toprule
\begin{minipage}[b]{0.24\columnwidth}\raggedright\strut
\textbf{CU-11}\strut
\end{minipage} & \begin{minipage}[b]{0.72\columnwidth}\raggedright\strut
\textbf{Envío de petición de cambio}\strut
\end{minipage}\tabularnewline
\cmidrule(r){1-2}
\textbf{Versión}       & 1.0           \\
\textbf{Autor}       & Adrián  Aguado    \\
\textbf{Requisitos asociados}       & RF-3.1 \\
\textbf{Descripción} & Gestión de los turnos \\
& entre usuarios. \\
\textbf{Precondición}  & El usuario está dentro de la \\
& aplicación.      \\
& Hay más de un usuario. \\
\textbf{Acciones} & 1. Enviar petición de cambio. \\
\textbf{Postcondición} & Ver si el usuario . \\
 & destinatario a aceptado o rechazado la petición. \\
\textbf{Excepciones} &   Aviso de envío de petición  \\
 &  tan solo una vez. \\
 &  Si no hay usuarios libres no es  \\
  &  posible enviar la petición  \\
\textbf{Frecuencia} & Media          \\
\textbf{Importancia} & Alta            \\
\textbf{Comentarios } & Se controlan que la       \\
& petición se envíe una sola vez. \\
& Se tienen en cuenta todos los usuarios      \\
& con turnos libre. \\
\bottomrule
\end{tabular}
\caption{CU-11 Enviar petición} 
\end{table}

%CU-12: Aceptar petición
\begin{table}[H]
\begin{tabular}{llr}  
\toprule
\begin{minipage}[b]{0.24\columnwidth}\raggedright\strut
\textbf{CU-12}\strut
\end{minipage} & \begin{minipage}[b]{0.72\columnwidth}\raggedright\strut
\textbf{Aceptar de petición de cambio}\strut
\end{minipage}\tabularnewline
\cmidrule(r){1-2}
\textbf{Versión}       & 1.0           \\
\textbf{Autor}       & Adrián  Aguado    \\
\textbf{Requisitos asociados}       & RF-3.1 \\
\textbf{Descripción} & Aceptación de la petición\\
\textbf{Precondición}  & El usuario está dentro de la \\
& aplicación.      \\
& Hay una petición por parte \\
& de otro usuario \\
\textbf{Acciones} & 1. Aceptar petición. \\
\textbf{Postcondición} & Cambio aceptado. \\
\textbf{Excepciones} &   Petición única  \\
\textbf{Frecuencia} & Media          \\
\textbf{Importancia} & Alta            \\
\textbf{Comentarios } & Se controlan que la       \\
& petición sea única. \\
\bottomrule
\end{tabular}
\caption{CU-12 Aceptar petición} 
\end{table}

%CU-13: Rechazar petición
\begin{table}[H]
\begin{tabular}{llr}  
\toprule
\begin{minipage}[b]{0.24\columnwidth}\raggedright\strut
\textbf{CU-13}\strut
\end{minipage} & \begin{minipage}[b]{0.72\columnwidth}\raggedright\strut
\textbf{Rechazar de petición de cambio}\strut
\end{minipage}\tabularnewline
\cmidrule(r){1-2}
\textbf{Versión}       & 1.0           \\
\textbf{Autor}       & Adrián  Aguado    \\
\textbf{Requisitos asociados}       & RF-3.2 \\
\textbf{Descripción} & Rechazo de la petición\\
\textbf{Precondición}  & El usuario está dentro de la \\
& aplicación.      \\
& Hay una petición por parte \\
& de otro usuario \\
\textbf{Acciones} & 1. Rechazar petición. \\
\textbf{Postcondición} & Fin de la lógica. \\
\textbf{Excepciones} &     \\
\textbf{Frecuencia} & Media          \\
\textbf{Importancia} & Alta            \\
\textbf{Comentarios } & Se controlan que la       \\
& petición sea única. \\
\bottomrule
\end{tabular}
\caption{CU-13 Rechazar petición} 
\end{table}

%CU-14: Aceptar intercambio
\begin{table}[H]
\begin{tabular}{llr}  
\toprule
\begin{minipage}[b]{0.24\columnwidth}\raggedright\strut
\textbf{CU-14}\strut
\end{minipage} & \begin{minipage}[b]{0.72\columnwidth}\raggedright\strut
\textbf{Aceptar intercambio}\strut
\end{minipage}\tabularnewline
\cmidrule(r){1-2}
\textbf{Versión}       & 1.0           \\
\textbf{Autor}       & Adrián  Aguado    \\
\textbf{Requisitos asociados}       & RF-3.3 \\
\textbf{Descripción} & Aceptación intercambio \\
\textbf{Precondición}  & El usuario está dentro de la \\
& aplicación.      \\
& Hay una petición por parte \\
& de otro usuario \\
\textbf{Acciones} & 1. Aceptar cambio. \\
\textbf{Postcondición} & Se produce el cambio. \\
& Los eventos se intercambian tanto \\
& en la base de datos como en el \\
& calendario de ambos usuarios \\
\textbf{Excepciones} &     \\
\textbf{Frecuencia} & Media          \\
\textbf{Importancia} & Alta            \\
\textbf{Comentarios } &   \\
\bottomrule
\end{tabular}
\caption{CU-14 Aceptar intercambio} 
\end{table}

%CU-15: Rechazar intercambio
\begin{table}[H]
\begin{tabular}{llr}  
\toprule
\begin{minipage}[b]{0.24\columnwidth}\raggedright\strut
\textbf{CU-15}\strut
\end{minipage} & \begin{minipage}[b]{0.72\columnwidth}\raggedright\strut
\textbf{Rechazar de petición de cambio}\strut
\end{minipage}\tabularnewline
\cmidrule(r){1-2}
\textbf{Versión}       & 1.0           \\
\textbf{Autor}       & Adrián  Aguado    \\
\textbf{Requisitos asociados}       & RF-3.4 \\
\textbf{Descripción} & Rechazo de la petición\\
\textbf{Precondición}  & El usuario está dentro de la \\
& aplicación.      \\
& Hay una petición por parte \\
& de otro usuario \\
\textbf{Acciones} & 1. Rechazar intercambio. \\
\textbf{Postcondición} & Fin de la lógica. \\
\textbf{Excepciones} &     \\
\textbf{Frecuencia} & Media          \\
\textbf{Importancia} & Alta            \\
\textbf{Comentarios } &  \\
\bottomrule
\end{tabular}
\caption{CU-15: Rechazar intercambio} 
\end{table}

%CU-16: Visualizar detalles 
\begin{table}[H]
\begin{tabular}{llr}  
\toprule
\begin{minipage}[b]{0.24\columnwidth}\raggedright\strut
\textbf{CU-16}\strut
\end{minipage} & \begin{minipage}[b]{0.72\columnwidth}\raggedright\strut
\textbf{Detalles peticiones}\strut
\end{minipage}\tabularnewline
\cmidrule(r){1-2}
\textbf{Versión}       & 1.0           \\
\textbf{Autor}       & Adrián  Aguado    \\
\textbf{Requisitos asociados}       & RF-3.6, RF-3.6.1  \\
& RF-3.6.2, RF-3.6.3, RF-3.6.4   \\
\textbf{Descripción} & Detalles de las peticiones\\
& mediante botones para una mejor \\
& comprensión por parte del usuario\\
\textbf{Precondición}  & El usuario está dentro de la \\
& aplicación.      \\
\textbf{Acciones} & 1. Visualizar normales. \\
& 2. Visualizar pendientes. \\
& 3. Visualizar aceptados. \\
& 4. Visualizar rechazados. \\
\textbf{Postcondición} & Ver numero de cada \\
& parámetro \\
\textbf{Excepciones} &     \\
\textbf{Frecuencia} & Media          \\
\textbf{Importancia} & Media           \\
\textbf{Comentarios } &  No es estrictamente necesaria \\
&  esta parte pero ayuda al usuario \\
\bottomrule
\end{tabular}
\caption{CU-16: Visualizar detalles} 
\end{table}

%CU-17: Visualizar normales
\begin{table}[H]
\begin{tabular}{llr}  
\toprule
\begin{minipage}[b]{0.24\columnwidth}\raggedright\strut
\textbf{CU-17}\strut
\end{minipage} & \begin{minipage}[b]{0.72\columnwidth}\raggedright\strut
\textbf{Detalles turnos}\strut
\end{minipage}\tabularnewline
\cmidrule(r){1-2}
\textbf{Versión}       & 1.0           \\
\textbf{Autor}       & Adrián  Aguado    \\
\textbf{Requisitos asociados}       & RF-3.6.1  \\
\textbf{Descripción} & Detalle de turnos \\
\textbf{Precondición}  & El usuario está dentro de la \\
& aplicación y ha realizado algún turno   \\
\textbf{Acciones} & 1. Visualizar normales. \\
\textbf{Postcondición} &  \\
\textbf{Excepciones} &     \\
\textbf{Frecuencia} & Media          \\
\textbf{Importancia} & Media           \\
\textbf{Comentarios } &   \\
\bottomrule
\end{tabular}
\caption{CU-17: Visualizar normales} 
\end{table}

%CU-18: Visualizar pendientes
\begin{table}[H]
\begin{tabular}{llr}  
\toprule
\begin{minipage}[b]{0.24\columnwidth}\raggedright\strut
\textbf{CU-18}\strut
\end{minipage} & \begin{minipage}[b]{0.72\columnwidth}\raggedright\strut
\textbf{Detalles peticiones pendientes}\strut
\end{minipage}\tabularnewline
\cmidrule(r){1-2}
\textbf{Versión}       & 1.0           \\
\textbf{Autor}       & Adrián  Aguado    \\
\textbf{Requisitos asociados}       & RF-3.6.2  \\
\textbf{Descripción} & Detalle de peticiones \\
& pendientes \\
\textbf{Precondición}  & El usuario está dentro de la \\
& aplicación y ha solicitado algún   \\
& cambio de turno (petición pendiente)   \\
\textbf{Acciones} & 1. Visualizar peticiones pendientes. \\
\textbf{Postcondición} &  \\
\textbf{Excepciones} &     \\
\textbf{Frecuencia} & Media          \\
\textbf{Importancia} & Media           \\
\textbf{Comentarios } &   \\
\bottomrule
\end{tabular}
\caption{CU-18 Visualizar pendientes} 
\end{table}

%CU-19: Visualizar aceptadas
\begin{table}[H]
\begin{tabular}{llr}  
\toprule
\begin{minipage}[b]{0.24\columnwidth}\raggedright\strut
\textbf{CU-19}\strut
\end{minipage} & \begin{minipage}[b]{0.72\columnwidth}\raggedright\strut
\textbf{Visualizar aceptadas}\strut
\end{minipage}\tabularnewline
\cmidrule(r){1-2}
\textbf{Versión}       & 1.0           \\
\textbf{Autor}       & Adrián  Aguado    \\
\textbf{Requisitos asociados}       & RF-3.6.3  \\
\textbf{Descripción} & Detalle de peticiones \\
& aceptadas\\
\textbf{Precondición}  & El usuario está dentro de la \\
& aplicación y ha realizado alguna  \\
& petición de cambio.  \\
\textbf{Acciones} & 1. Visualizar peticiones\\
& aceptadas . \\
\textbf{Postcondición} &  Confirmar cambio \\
\textbf{Excepciones} &  Si la petición se ha rechazado   \\
&  no se visualizará aquí.   \\
\textbf{Frecuencia} & Media          \\
\textbf{Importancia} & Media           \\
\textbf{Comentarios } &   \\
\bottomrule
\end{tabular}
\caption{CU-19: Visualizar aceptadas} 
\end{table}

%CU-20: Visualizar rechazadas
\begin{table}[H]
\begin{tabular}{llr}  
\toprule
\begin{minipage}[b]{0.24\columnwidth}\raggedright\strut
\textbf{CU-20}\strut
\end{minipage} & \begin{minipage}[b]{0.72\columnwidth}\raggedright\strut
\textbf{Detalles peticiones}\strut
\end{minipage}\tabularnewline
\cmidrule(r){1-2}
\textbf{Versión}       & 1.0           \\
\textbf{Autor}       & Adrián  Aguado    \\
\textbf{Requisitos asociados}       & RF-3.6.4  \\
\textbf{Descripción} & Detalle de peticiones \\
& rechazadas \\
\textbf{Precondición}  & El usuario está dentro de la \\
& aplicación y ha realizado alguna  \\
& petición de cambio.  \\
\textbf{Acciones} & 1. Visualizar peticiones\\
& rechazadas . \\
\textbf{Postcondición} &  Rechazar cambio \\
\textbf{Excepciones} &  Si la petición se ha aceptado   \\
&  no se visualizará aquí.   \\
\textbf{Frecuencia} & Media          \\
\textbf{Importancia} & Media           \\
\textbf{Comentarios } &   \\
\bottomrule
\end{tabular}
\caption{CU-20 Visualizar rechazadas} 
\end{table}

%CU-21: Visualizar Calendario
\begin{table}[H]
\begin{tabular}{llr}  
\toprule
\begin{minipage}[b]{0.24\columnwidth}\raggedright\strut
\textbf{CU-21}\strut
\end{minipage} & \begin{minipage}[b]{0.72\columnwidth}\raggedright\strut
\textbf{Visualizar calendario}\strut
\end{minipage}\tabularnewline
\cmidrule(r){1-2}
\textbf{Versión}       & 1.0           \\
\textbf{Autor}       & Adrián  Aguado    \\
\textbf{Requisitos asociados}       & RF-3.6.4  \\
\textbf{Descripción} & Visualizar cambios en \\
& el  calendario \\
\textbf{Precondición}  & El usuario está dentro de la \\
& aplicación. \\
\textbf{Acciones} & 1. Visualizar turnos \\
& 2. Click encima de turnos para ver \\
&  duración  \\
&  3, Color asociado a cada tipo de turno  \\
\textbf{Postcondición} & \\
\textbf{Frecuencia} & Media          \\
\textbf{Importancia} & Media           \\
\textbf{Comentarios } &  Todos los cambios de turno  \\
 &  se deben de ver reflejados en el calendario  \\
  &  del usuario en caso de ser aceptados por  \\
    &  ambas partes \\
\bottomrule
\end{tabular}
\caption{CU-21 Visualizar calendario} 
\end{table}






\apendice{Especificación de diseño}
\section{Introducción}\label{introducciondiseno}
En este apartado se incluye el diseño de la aplicación y se mostrarán también los diagramas de clases de cada una de las partes que forman el proyecto.

\section{Diseño de datos}\label{datos}

En este apartado se incluye el diseño de la aplicación y se mostrarán también los diagramas de clases de cada una de las partes que forman el proyecto.






\apendice{Documentación técnica de programación}

\section{Introducción}\label{introduccion-programador}
En este capítulo vamos a dividirlo en dos partes, por un lado los condiciones específicas propias de la aplicación y por otro vamos a aprender de manera breve y sencilla a crear una aplicación gracias a \emph{Angular CLI}, la interfaz de línea de comandos de Angular.

El por qué he decidido incluir la parte de \emph{Angular CLI} es por que considero que va aclarar muchas dudas acerca de cómo funciona esta tecnología y también para intentar comprender mejor como he realizado el proyecto. En diversas ocasiones para aprender a usar una tecnología tan solo hace falta tiempo pero en otras muchas el tiempo es limitado por lo que quizás con esta pequeña guía podemos comenzar a utilizar un framework cliente en muy pocos pasos y de una manera sencilla. Si bien es cierto que yo he necesitado mucho tiempo para comprender su funcionamiento y no lo he aprendido todo gracias a la interfaz de comandos, ésta te permite agilizar los trámites de creación en un tanto por ciento considerablemente alto.

Este anexo por tanto tiene como objetivo analizar y documentar las necesidades funcionales que deberán ser soportadas por el sistema a desarrollar, es decir, en qué condiciones ha sido desarrollado, en qué condiciones se debe usar y cuáles son los requerimientos mínimos para que un futuro programador interactúe con la aplicación en caso de que así lo considere.

\section{Requerimientos mínimos necesarios}
Además de un \emph{IDE} para poder trabajar, los prerequisitos mínimos para empezar a trabajar en el proyecto son:

\begin{enumerate}
	\item Instalar \hyperlink{https://nodejs.org/en/}{Nodejs} y \hyperlink{https://www.mongodb.com/}{MongoDB}.
	\item Instalar \hyperlink{https://www.npmjs.com/}{npm}
	\item Instalar Angular CLI
		\lstset{language=C, breaklines=true, basicstyle=\footnotesize}
		\begin{lstlisting}[frame=single]
		npm i -g @angular/cli
    	\end{lstlisting}
\end{enumerate}

Nota: es interesante, una vez instalados asegurarnos de que están correctamente instalados usando los comandos \emph{node -v} y \emph{npm -v} para saber si tenemos los requerimientos correctamente instalados.


\section{AngularCLI}\label{angularCLI}
Sabiendo los requisitos mínimos que son necesarios para proceder vamos a realizar, cómo he nombrado antes una breve introducción a la línea de comandos de Angular,  \textbf{Angular CLI}.

\subsection{¿Qué es typescript?}\label{typesript}
Esto se ha nombrado antes pero resulta esencial para entender el código de Angular. Antes de nada decir que en Angular2 se puede trabajar con Javascript “clásico” (ES5), así como ES6 y TypeScript. Desarrollado por Microsoft y como apuesta de Google vamos a ver el crecimiento exponencial que ha tenido este de este lenguaje, que no deja de ser un super conjunto de JavaScript. 
 TypeScript es por tanto un superset de ECMAScript 6, es decir, incluye todas las funcionalidades de ES6, y además incorpora una capa por encima con funcionalidades extra.

\begin{center}
    \begin{minipage}{0.9\linewidth}
        \vspace{5pt}%margen superior de minipage
        {\small
            \emph{TypeScript is a typed superset of JavaScript that compiles to plain JavaScript. Any browser. Any host. Any OS. Open source.}
        }
        \begin{flushright}
            (\cite{TypeScript Web})
        \end{flushright}
        \vspace{5pt}%margen inferior de la minipage
    \end{minipage}
\end{center}

La principal característica de TypeScript por encima de Javascript es que permite definir de qué tipo son las variables que se van a usar.

\subsection{Comandos básicos}\label{cbasicos}
\textbf{Angular CLI} nos va a hacer la vida más fácil con Angular, ya que nos permite crear de forma sencilla un aplicación lista para funcionar después de la instalación por línea de comandos con \emph{npm}, además incorpora las buenas prácticas de programación con Angular. 

\imagen{angularclicommands}{Comandos  \emph{AngularCLI}. Fuente:  \url{https://cli.angular.io/}.}

Por lo tanto para crear nuestra aplicación con esta herramienta tan solo debemos ir a nuestra carpeta donde queremos crear el proyecto e introducir en la línea de comandos: 
	\lstset{language=C, breaklines=true, basicstyle=\footnotesize}
		\begin{lstlisting}[frame=single]
		ng new angular-cli-primer-proyecto
    	\end{lstlisting}
 
 Es un proceso que durará unos 2 minutos,  Al finalizar seremos capaces de ejecutar directamente la aplicación por defecto que crea la herramienta. Para ejecutarla, al igual que antes accedemos a la carpeta del proyecto y ejecutamos:
 	\lstset{language=C, breaklines=true, basicstyle=\footnotesize}
		\begin{lstlisting}[frame=single]
		 npm start
    	\end{lstlisting}
    	
    	\imagen{angularcliconsola}{Consola ejecución. Fuente:  \url{http://codigoxules.org/}.}
    	
 Si te fijas en la ejecución se llama al comando \textbf{ng serve} que será el que utilicemos con Angular CLI para trabajar. Resulta realmente útil al introducir el último comando nombrado ya que una vez lanzado se detectan automáticamente los cambios a medida que vamos avanzando la página se recarga al guardar el componente que estemos modificando.  
    	
    	\imagen{angularcliresultado}{Resultado después de creado. Fuente:  \url{http://codigoxules.org/}.}
    	
   Esta es la sencilla manera de crear un proyecto con \textbf{Angular CLI}. Después existen una serie de comandos para crear directamente componentes, servicios, clases, interfaces.. etcétera.  Algunos de ellos son:
   
   
\begin{table}[htbp]
\begin{center}
\begin{tabular}{|l|l|}
\hline
Tipo a crear & Comando  \\
\hline \hline
Component & \emph{ng component my-new-component}  \\ \hline
Directive & \emph{ng g directive my-new-directive}  \\ \hline
Pipe & \emph{ng g pipe my-new-pipe}  \\ \hline
Service & \emph{ng g service my-new-service}  \\ \hline
Class & \emph{ng g class my-new-class}  \\ \hline
Interface & \emph{ng g interface my-new-interface}  \\ \hline
Enum & \emph{ng g enum my-new-enum}  \\ \hline
\end{tabular}
\caption{Tabla comandos Angular CLI}
\label{tabla:sencilla}
\end{center}
\end{table}
   


\section{Manual del programador}
En esta sección hay que tener en cuenta que el autor de este trabajo proyecto a escogido unas serie de herramientas, tanto para desplegar la app, como la base de datos como para desarrollar la aplicación pero que de ninguna manera resultan ser ni las únicas ni las mejores simplemente son unas herramientas que ha considerado utilizar pero existen muchas más que no son ni peores ni mejores.

\subsection{Prerequisitos}
\begin{enumerate}
	\item Instalar \hyperlink{https://nodejs.org/en/}{Nodejs} y \hyperlink{https://www.mongodb.com/}{MongoDB}.
	\item Instalar \hyperlink{https://www.npmjs.com/}{npm}
	\item Instalar Angular CLI
		\lstset{language=C, breaklines=true, basicstyle=\footnotesize}
		\begin{lstlisting}[frame=single]
		npm i -g @angular/cli
    	\end{lstlisting}
    \item Descargar el proyecto ( \emph{git clone} ) o bien desde un soporte
	\item Una vez tenemos el proyecto descargado en nuestro sistema local desde el directorio raíz (vía línea de comandos) instalar las dependencias necesarias.
	\lstset{language=C, breaklines=true, basicstyle=\footnotesize}
		\begin{lstlisting}[frame=single]
		npm install
    	\end{lstlisting}
\end{enumerate}


\subsection{Estructura de directorios }

La estructura de directorios general es la siguiente:

%Foto erronea cambiar!
\imagen{directorios}{Vista carpetas general  \emph{BarterAPP}. Fuente: Elaboración propia.}

\begin{itemize}
	\item \textbf{$/client/$}: \emph{Front-End} 
	\begin{itemize}
	\item \textbf{$/client/app$}: Carpeta App
		\begin{itemize}
		\item \textbf{$/client/app/_directives$}: hola
		\item \textbf{$/client/app/_guards$}: hola
		\item \textbf{$/client/app/_models$}: hola
		\item \textbf{$/client/app/_services$}:  hola
		\item \textbf{$/client/app/enviroments$}:  hola
		\item \textbf{$/client/app/pages/$}:  
			\begin{itemize}
			\item \textbf{$/client/app/pages/angular_calendar $}:  hola
			\item \textbf{$/client/app/pages/home $}:   hola
			\item \textbf{$/client/app/pages/login $}:   hola
			\item \textbf{$/client/app/pages/register $}:   hola
			\item \textbf{$/client/app/pages/welcome $}:   hola
			\end{itemize}
		\item \textbf{$/client/app/repository $}:  
		\item \textbf{$/client/app/utils $}:  
		\end{itemize}
	\item \textbf{$/client/assets $}: Recursos  
	\item \textbf{$/client/enviroments $}:  Entornos
   \end{itemize}
	\item \textbf{$/dist/ $}:
	\item \textbf{$/e2e/ $}:
	\item \textbf{$/node_modules/ $}:
	\item \textbf{$/server/ $}:  \emph{Back-End} 
		\begin{itemize}
			\item \textbf{$/server/config $}:
			\item \textbf{$/server/controllers $}:
			\item \textbf{$/node_modules/ $}:
			\item \textbf{$/server/routes $}:
			\item \textbf{$/server/services $}:
		\end{itemize}
\end{itemize}

Observables
http://academia-binaria.com/comunicaciones-http-observables-con-angular2/

angular2 polling -observables
https://stackoverflow.com/questions/41658162/how-to-do-polling-with-angular-2-observables

Guards
https://blog.thoughtram.io/angular/2016/07/18/guards-in-angular-2.html

Interpolacion
http://academia-binaria.com/databinding-el-flujo-de-datos-de-angular2/

DTO
https://stackoverflow.com/questions/39272947/dto-design-in-typescript-angular2

ngAfterViewInit
https://angular.io/docs/ts/latest/api/core/index/AfterViewInit-class.html


PrimeNg
%https://www.primefaces.org/primeng/#/

Growl
%https://www.primefaces.org/primeng/#/growl

Login/Register tutorial
%http://jasonwatmore.com/post/2016/09/29/angular-2-user-registration-and-login-example-tutorial#register-component-html

AngularCalendar
https://github.com/mattlewis92/angular-bootstrap-calendar

Multiple components
https://angular.io/docs/ts/latest/tutorial/toh-pt3.html

Tutorial angular2 components
https://scotch.io/tutorials/creating-your-first-angular-2-components

https://desarrolloweb.com/articulos/analisis-carpetas-proyecto-angular2.html


https://desarrolloweb.com/articulos/introduccion-componentes-angular2.html

http://academia-binaria.com/hola-mundo-en-angular-2/


\subsection{Modo desarrollador}

Hay que tener en cuenta que tenemos diversas tecnologías dentro del proyecto por lo que resulta un poco costoso tener que ejecutar todas por separado, que es lo que hacía en un principio. Con el modo desarrollador (emula al \textbf{ng serve} nombrado antes, tan solo con el comando siguiente ejecutamos MongoDB, Angular, Express y el compilador de Typescript luego ya no nos hace falta nada más y podemos comenzar a trabajar. 

	\lstset{language=C, breaklines=true, basicstyle=\footnotesize}
		\begin{lstlisting}[frame=single]
		npm run dev
    	\end{lstlisting}

\subsection{Servidor (Heroku)}
Para subir la aplicación al servidor yo he elegido \hyperlink{https://www.heroku.com/}{Heroku} y los pasos para lanzar la aplicación son los que se enumeran a continuación:

\begin{enumerate}
	\item Acceder a la \hyperlink{https://www.heroku.com/}{web} de \emph{Heroku}, registrarse y crear una nueva app.
	\item Instalar \hyperlink{https://devcenter.heroku.com/articles/heroku-cli}{Heroku CLI}
	\item Ejecutar los siguientes comandos:
		\lstset{language=C, breaklines=true, basicstyle=\footnotesize}
		\begin{lstlisting}[frame=single]
		heroku login
		cd my-project/
		git init
		heroku git:remote -a your-app-name
    	\end{lstlisting}
	\item Descargar el proyecto BarterApp desde github 
	\item Editar el servidor(Ruta: \emph{server/config/db.ts}) para introducir un servidor real con MongoDB (lo veremos en la sección siguiente)
	\item Ejecutar los siguientes comandos:
		\lstset{language=C, breaklines=true, basicstyle=\footnotesize}
		\begin{lstlisting}[frame=single]
		npm i
		ng build -prod or ng build -aot -prod
		tsc -p server
		git add .
		git commit -m "Going to Heroku"
		git push heroku master
		heroku open
    	\end{lstlisting}
	\item Una venta se abrirá con tu aplicación ya online
\end{enumerate}

Resumiendo de esta manera utilizamos Heroku como si de github se tratase y vamos aplicando los cambios que hacemos a nuestra aplicación al repositorio creado en Heroku, la ventaja que ya tenemos la aplicación lanzada en un entorno real. Es una restricción importante que ya no podemos trabajar con una base de datos local sino que tiene que ser un servidor remoto con MongoDB ( en este caso) para que la aplicación funcione.
	

\subsection{Base de datos (Mlab)}
La base de datos escogida es MongoDB, la instalación es fácil y sencilla tan solo hay que ir a la \hyperlink{http://www.mongodb.org/downloads}{web} oficial y descargar la última versión estable disponible. Una vez en nuestro ordenador la ejecutamos y listo.  Si tenéis algún problema os recomiendo seguir alguno de los tutoriales que existen en la web. https://www.adictosaltrabajo.com/tutoriales/mongodb/

Como sabemos podemos manejar la base de datos desde línea de comandos pero resulta un poco pesado cuando empezamos a tener muchas colecciones o un número de datos alto. Es por eso que yo recomiendo encarecidamente  \hyperlink{https://robomongo.org/}{Robomongo}, la copia de SQL server para bases de datos NoSQL.  Interfaz no muy lograda pero perfecta para visualizar nuestras bases de datos e información de una manera más intuitiva.

	\imagen{robomongoview}{Vista Robomongo. Fuente:  Elaboración propia.}
	
	Esta claro que cuando trabajamos con un proyecto en local para realizar pruebas lo mejor es tener una base de datos en local pero cuando queremos algo un poco más serio resulta interesante el tener esa base de datos en un servidor real. Para ello he escogido \hyperlink{mlab.com}{Mlab}, los pasos para crear la base de datos se enumeran a continuación: 
	
	\begin{enumerate}
	\item Acceder a la \hyperlink{mlab.com}{Mlab}, registrarse y crear una cuenta
	\item Crear una subscripción a una base de datos (Tiene de pago pero también tiene una versión gratuita que nos vale para realizar pruebas a este nivel)
	\item Ejecutar el siguiente comando para establecer la conexión con la base de datos creada
		\lstset{language=C, breaklines=true, basicstyle=\footnotesize}
		\begin{lstlisting}[frame=single]
		mongo ds012345.mlab.com:56789/dbname -u dbuser -p dbpassword
    	\end{lstlisting}
	\item Las colecciones las podemos crear directamente desde la propia web, o bien desde la lína de comandos 
	\item Para probar que la base de datos funciona (El primer comando es para insertar un elemento y el segundo para buscar. En este caso el segundo debería devolver el único elemento insertado.:
		\lstset{language=C, breaklines=true, basicstyle=\footnotesize}
		\begin{lstlisting}[frame=single]
		db.mynewcollection.insert({ "foo" : "bar" })
		db.mynewcollection.find() 
    	\end{lstlisting}	
\end{enumerate}

Nota: para empezar a trabajar con este proyecto es necesario crear las tres colecciones de las que ya hemos hablado en secciones anteriores: \emph{Users, Events, Interchange}, y \underline{no} hacer nada más puesto que según la configuración de la aplicación se debería añadir los usuarios, eventos y turnos de manera automática cuando el usuario hacer click en los botones adecuados. 


\begin{itemize}
  \item Stack MEAN
  \begin{itemize}
    \item \textbf{M}ongoose.js (MongoDB):  base de datos
    \item \textbf{E}xpress.js: backend framework
    \item \textbf{A}ngular 4: frontend framework
    \item \textbf{N}ode.js: entorno tiempo de ejecución
  \end{itemize}
  \item Otras tecnologías que se utilizan:
  \begin{itemize}
    \item Angular CLI
    \item Boostrap
    \item Font Awesome
    \item JSON WEB TOKEN
  \end{itemize}
\end{itemize}





\apendice{Documentación de usuario}

\section{Introducción}\label{introduccion-usuario}
En este capítulo se detalla como un usuario puede comenzar a usar barterAPP.

\section{Instalación}\label{instalacion-app}
\subsection{Descarga}\label{descarga}

\section{Manual del usuario}\label{manual-usuario}




\bibliography{bibliografiaAnexos}
\bibliographystyle{plainnat}








\end{document}